\section{魏的朝廷与政治社会：名位怎样变成可用的秩序}

魏的洛阳并不只生产皇帝的命令，也生产让命令看似合乎次序的官名、礼仪、学校
和品评。受禅以后，旧汉官员、新近依附的军人和地方名士需要共同使用一套能够
任命、奏事与处分的语言。曹丕定都洛阳的价值正在这里：尚书台、宫卫和通向
河北、关中、淮南的道路可以在近畿相接。它并没有消灭地方距离；反而使每次
边地危机都显示中枢的文书究竟能否换成仓粮、援军与可信的使者。

九品官人法是这种接合的一个入口。陈群的制度常被后世概括为门阀政治的开端，
但在黄初的环境中，它首先处理的是新朝如何取得可任用的人。乡里评价、州郡
官员与中央选任被放进等级次序，减少了每次任官都从零开始辨认的成本。代价是
谁有资格作保、谁的家世和迁徙经历可被看见，会在进入洛阳以前就影响名位。
《三国志·陈群传》能说明制度与人物的关系，不能提供全国品评的完整名册；
不能把这项安排写成一夜之间固定全社会的铁律。

石经残片补足了另一种中枢景象。正始三体石经被刻置于太学，字形与经文因而
能够被公开校读。残石让我们看见刻工、阅读者和官学空间，而不是抽象的“儒学
复兴”。它也不证明各郡都读同一种字或同一种经。品第关心谁可进入官署，石经
关心何种文本能够被共同援引；两者共同构成魏朝尝试制造政治社会的方式，却都
要依赖地方人际网络才能真正发挥作用。

曹叡以后，辅政危机改变了这些制度的重心。曹爽与司马懿争夺的不是一件名为
权力的器物，而是宿卫、诏令、任官与中军的接头。高平陵之变中，城门、太后
名义和官署文书必须同时被一方接住，才足以排除另一方。传记中的承诺与指控多
在胜负已定后被保存，不能充当逐刻记录；但它们提示读者，政治秩序的脆弱不在
于缺少礼仪，而在于礼仪所指向的军政入口已可被少数人抢先控制。

淮南的连续叛乱把洛阳的这一裂缝带到水陆走廊。王凌、毌丘俭、文钦和诸葛诞
各有不同的同盟、城池和继承主张，不能被化成同一类“叛臣”。寿春的粮道、吴
援能否穿过淮河、魏廷是否能先送来赦令与援军，决定他们的选择何时还可维持。
司马氏最终能够压下这些危机，说明它掌握的不是抽象声望，而是一条比对手更能
运转的命令链。也正因为如此，魏的名义可以继续用于日常政务，直到司马氏认为
公开受禅的成本已经较低。
